% ==============================================================================
% ARSIP REFERENSI / DAFTAR PUSTAKA CADANGAN
% Proyek: Proposal Tugas Akhir - Sistem Pencegahan Diskrepansi POS
% Dokumen ini menyimpan referensi tambahan yang sewaktu-waktu dapat disalin
% kembali ke file utama (proposal.tex) apabila diperlukan oleh dosen pembimbing.
% ==============================================================================

% ------------------------------------------------------------------------------
% KATEGORI 1: POINT OF SALES (POS) & SISTEM KASIR
% ------------------------------------------------------------------------------

% 1. Angelline & Dolok (2026) - Rancang bangun POS berbasis Web
\noindent Angelline, C., \& Dolok, M. (2026). Perancangan Aplikasi Point of Sales Berbasis Web pada Toko Niaga. \textit{Jurnal Ilmu Komputer dan Sistem Informasi}, 14(1).

\vspace{0.2cm}
% 2. Ghozali & Iskendang (2020) - Aplikasi POS Web untuk Koperasi
\noindent Ghozali, A. L., \& Iskendang. (2020). Penerapan Aplikasi Point of Sales (POS) untuk Menentukan Jasa Anggota dalam Transaksi Penjualan pada Koperasi Berbasis Web. \textit{Jurnal Ilmiah Ilmu Komputer}, 6(1), 27--30.

\vspace{0.2cm}
% 3. Kurniawan (2026) - REST API & Dashboard Admin POS Multistore (Laravel)
\noindent Kurniawan, B. D. (2026). \textit{Perancangan dan Implementasi REST API serta Dashboard Admin pada Sistem Point of Sale (POS) FLOKASIR Multistore Menggunakan Laravel} [Laporan Kerja Praktek]. Politeknik Negeri Bengkalis.

\vspace{0.4cm}
% ------------------------------------------------------------------------------
% KATEGORI 2: REST API, WEB SERVICE & BACKEND FRAMEWORK
% ------------------------------------------------------------------------------

% 4. Gowell & Suprihadi (2024) - REST API Laravel untuk Web Service
\noindent Gowell, K., \& Suprihadi. (2024). Perancangan Web Service REST API Menggunakan PHP dan Framework Laravel di Tenta Tour Salatiga. \textit{Jurnal JTIK (Jurnal Teknologi Informasi dan Komunikasi)}, 8(1), 49--57. \url{https://doi.org/10.35870/jtik.v8i1.1269}

\vspace{0.2cm}
% 5. Herdiyatmoko (2022) - Desain Backend REST API Laravel
\noindent Herdiyatmoko, H. F. (2022). Desain Sistem Backend Berbasis Rest Api Menggunakan Framework Laravel 7. \textit{SKANIKA}, 5(2), 136--144. \url{https://doi.org/10.36080/skanika.v5i2.2947}

\vspace{0.2cm}
% 6. Abidin et al. (2024) - Pengembangan Web Laravel dengan Metode Waterfall
\noindent Abidin, Z., Kurniawan, J., Suntoro, D. A. P., \& Yuliastuti, G. E. (2024). Rancang Bangun Website Profil Ikatan Mahasiswa Arosbaya Menggunakan Laravel dan Bootstrap 5 Dengan Metode Waterfall. \textit{Prosiding Seminar Implementasi Teknologi Informasi dan Komunikasi}, 3(2), 257--268.

\vspace{0.4cm}
% ------------------------------------------------------------------------------
% KATEGORI 3: METODE ANALISIS ABC / INVENTORY PENGEMBANGAN LANJUTAN
% ------------------------------------------------------------------------------

% 7. Flores & Whybark (1986) - Fondasi Teori Multiple Criteria ABC Analysis
\noindent Flores, B. E., \& Whybark, D. C. (1986). Multiple criteria ABC analysis. \textit{International Journal of Operations \& Production Management}, 6(3), 38--46. \url{https://doi.org/10.1108/eb054765}
